\section{Quantitative bounds and separation}
\label{sec:separation}

The dominant terms $M^2S'_0$ and $S'_1$ in
Proposition~\ref{prop:recurrences} suggest normalizing $S_{\ell,b}$ by
$M^{\ell+b}$. The joint profiles vanish initially; choosing $M$ large
relative to $d$ keeps their contribution small through the prescribed
depth. The following bound makes this choice explicit.

\begin{lemma}
\label{lem:quantitative-profiles}
Suppose that $d\geq 1$ and $M\geq 9^{d+1}$.
For every $0\leq\ell\leq d$ and $b\in\{0,1\}$, the predicates of the
construction satisfy
\[
  S_{\ell,b}\leq 2L M^{\ell+b}.
\]
\end{lemma}

\begin{proof}
Define
\[
 u_\ell=\max_{b\in\{0,1\}}\frac{S_{\ell,b}}{M^{\ell+b}},
 \qquad
 v_\ell=\max_{b\in\{0,1\}}\frac{J_{\ell,b}}{M^{\ell+b}},
 \qquad
 \epsilon=\frac{\max(r,t)}{M}.
\]
The base bounds in~\eqref{eq:base-profile} give $u_0\leq L$ and $v_0=0$.
For $1\leq\ell\leq d$, dividing the four inequalities in
Proposition~\ref{prop:recurrences} by the corresponding powers of~$M$
gives
\begin{align*}
 \frac{S_{\ell,1}}{M^{\ell+1}}
   &\leq \left(1+\frac rM\right)u_{\ell-1},
 &\qquad
 \frac{S_{\ell,0}}{M^\ell}
   &\leq \left(1+\frac tM\right)u_{\ell-1}+3v_{\ell-1},\\
 \frac{J_{\ell,1}}{M^{\ell+1}}
   &\leq \frac rM u_{\ell-1}+v_{\ell-1},
 &
 \frac{J_{\ell,0}}{M^\ell}
   &\leq \frac tM u_{\ell-1}+3v_{\ell-1}.
\end{align*}
Consequently,
\begin{equation}
\label{eq:normalized-recurrences}
 \begin{split}
  u_\ell&\leq (1+\epsilon)u_{\ell-1}+3v_{\ell-1},\\
  v_\ell&\leq \epsilon u_{\ell-1}+3v_{\ell-1}.
 \end{split}
\end{equation}
The forcing term for $v_\ell$ is at most $\epsilon u_{\ell-1}$, while
its inherited contribution is multiplied by at most three. We therefore
make $\epsilon$ small compared with $3^{-d}$. Since $M\geq81$, both
$r=\lceil\sqrt M\rceil\leq2\sqrt M$ and $t=16\leq2\sqrt M$.
It follows that
\begin{equation}
\label{eq:epsilon-choice}
 \epsilon\leq\frac{2}{\sqrt M}\leq\frac{2}{3^{d+1}}.
\end{equation}

Let $U_0=L$, $V_0=0$, and define the comparison sequences by
\[
 U_\ell=(1+\epsilon)U_{\ell-1}+3V_{\ell-1},
 \qquad
 V_\ell=\epsilon U_{\ell-1}+3V_{\ell-1}.
\]
All coefficients are nonnegative, so~\eqref{eq:normalized-recurrences}
implies $u_\ell\leq U_\ell$ and $v_\ell\leq V_\ell$ by induction.
We prove simultaneously that, for $0\leq\ell\leq d$,
\begin{equation}
\label{eq:comparison-bounds}
 U_\ell\leq2L,
 \qquad
 V_\ell\leq\epsilon L(3^\ell-1).
\end{equation}
These bounds hold at $\ell=0$.
Suppose they hold at every index less than some $1\leq\ell\leq d$.
First,
\[
 V_\ell
 \leq 2\epsilon L+3\epsilon L(3^{\ell-1}-1)
 =\epsilon L(3^\ell-1).
\]
Next, summing the increments of $U_j$ and using only the bounds at
indices $j<\ell$, we obtain
\begin{align*}
 U_\ell
 &=L+\sum_{j=0}^{\ell-1}(\epsilon U_j+3V_j)\\
 &\leq L+2\epsilon L\ell
          +3\epsilon L\sum_{j=0}^{\ell-1}(3^j-1)\\
 &=L+\epsilon L\left[\frac32(3^\ell-1)-\ell\right]
 <2L.
\end{align*}
The last inequality follows from~\eqref{eq:epsilon-choice}, since
\[
 \epsilon\left[\frac32(3^\ell-1)-\ell\right]
 \leq\frac{3^{\ell+1}-3-2\ell}{3^{d+1}}<1.
\]
This completes the simultaneous induction. In particular,
$u_\ell\leq2L$, which is the desired bound.
\end{proof}

\begin{proof}[Proof of Theorem~\ref{thm:main}]
Fix any integer $d\geq1$ and choose an integer $M\geq9^{d+1}$.
Lemma~\ref{lem:labelling} provides the required tournament labels.
Following the OR-balancing principle of Ambainis and Sun
\cite[Section~4, Lemma~1]{AmbainisSun2011}, take $f$ to be the OR of $M$
disjoint copies of $P_{d,1}$.
Each copy has dimension $LM(kr)^d$, so $f$ is a total Boolean function on
\begin{equation}
\label{eq:final-dimension}
 n=LM^2(kr)^d
\end{equation}
bits. By Lemma~\ref{lem:blocks}, every copy is zero at its zero input,
and a block from its stated family makes that copy one.
Thus $f(0)=0$ and $f$ is nonconstant.

At a zero input of $f$, every copy is zero, and its sensitivity is at most
$M S_{d,0}$. At a one input, either at least two copies accept, in which
case no single bit can change the OR to zero, or exactly one copy accepts.
In the latter case every sensitive bit belongs to that accepting copy,
so there are at most $S_{d,1}$ such bits.
Lemma~\ref{lem:quantitative-profiles} therefore gives
\begin{equation}
\label{eq:final-sensitivity}
 s(f)\leq\max\{M S_{d,0},S_{d,1}\}
       \leq2L M^{d+1}.
\end{equation}

Each copy supplies $M k^d$ pairwise disjoint blocks sensitive at zero,
by Lemma~\ref{lem:blocks}. Taking their union over the disjoint copies
gives a family of $M^2k^d$ pairwise disjoint sensitive blocks for $f$ at zero.
In particular,
\begin{equation}
\label{eq:final-block-sensitivity}
 \bs(f,0)\geq M^2k^d.
\end{equation}
Since $f$ is nonconstant, $s(f)>0$. Combining
\eqref{eq:final-sensitivity} and~\eqref{eq:final-block-sensitivity} yields
\begin{equation}
\label{eq:ratio-separation}
 \frac{\bs(f)}{s(f)^2}
 \geq\frac{\bs(f,0)}{s(f)^2}
 \geq\frac{(k/M^2)^d}{4L^2}
 \geq\frac{2^d}{4(d+2)^2}.
\end{equation}
The last expression tends to infinity as $d$ tends to infinity.
For any proposed constant $C$, first choose $d$ so that this expression
exceeds $C$, then choose $M\geq9^{d+1}$, then choose labels supplied by
Lemma~\ref{lem:labelling} and form $f$ as above.
The resulting finite-dimensional total nonconstant function satisfies
$\bs(f)>C s(f)^2$. Hence no universal quadratic constant exists.
\end{proof}

\subsection{A fixed exponent by composition}

The construction also supplies a fixed seed whose block sensitivity at
zero exceeds the square of its sensitivity. Repeated composition of that
seed gives a power separation. This is the composition-and-powering
principle studied by Tal~\cite{Tal2013}. We include the elementary composition facts
to make this additional consequence self-contained.

\begin{lemma}
\label{lem:composition}
Let $f:\{0,1\}^a\to\{0,1\}$ and
$g:\{0,1\}^b\to\{0,1\}$ be total Boolean functions.
Their disjoint composition is the function on $ab$ bits defined by
\[
 (f\circ g)(x_1,\ldots,x_a)
   =f\bigl(g(x_1),\ldots,g(x_a)\bigr),
 \qquad x_i\in\{0,1\}^b.
\]
Then
\[
 s(f\circ g)\leq s(f)s(g).
\]
If $g(0)=0$, then $(f\circ g)(0)=f(0)$ and
\[
 \bs(f\circ g,0)\geq\bs(f,0)\bs(g,0).
\]
More explicitly, any disjoint families of $p$ and $q$ blocks sensitive at
zero for $f$ and $g$, respectively, yield a disjoint family of $pq$ blocks
sensitive at zero for $f\circ g$.
\end{lemma}

\begin{proof}
For an input $x=(x_1,\ldots,x_a)$, put
$y=(g(x_1),\ldots,g(x_a))$.
A bit in child $x_i$ is sensitive for $f\circ g$ precisely when it is
sensitive for $g$ at $x_i$ and coordinate $i$ is sensitive for $f$ at $y$.
Thus
\[
 s(f\circ g,x)
   =\sum_{i:\,f(y^{\{i\}})\ne f(y)}s(g,x_i)
   \leq s(f,y)s(g)
   \leq s(f)s(g).
\]

Now suppose $g(0)=0$. Let $C_1,\ldots,C_p\subseteq[a]$ and
$D_1,\ldots,D_q\subseteq[b]$ be disjoint nonempty sensitive blocks for
$f$ and $g$ at zero. Identify the coordinates of the composition with
$[a]\times[b]$. For each $1\leq i\leq p$ and $1\leq j\leq q$, set
\[
 E_{i,j}=C_i\times D_j.
\]
These blocks are nonempty and pairwise disjoint: different $i$ use
disjoint sets of children, while for fixed $i$ different $j$ use
disjoint coordinate sets in each child.
Flipping $E_{i,j}$ at zero makes precisely the children indexed by $C_i$
output one, because $g(0)=0$ and $D_j$ is sensitive there.
The outer input therefore becomes $0^{C_i}$, which changes the value of
$f$. Hence every $E_{i,j}$ is sensitive for $f\circ g$ at zero.
Choosing maximum families proves the stated block-sensitivity inequality;
if either maximum is zero, the inequality is immediate.
\end{proof}

\begin{corollary}
\label{cor:power-separation}
There exist a real number $\alpha>2$ and nonconstant total Boolean
functions $F_m$, for integers $m\geq1$, such that
\[
 \bs(F_m,0)\geq s(F_m)^\alpha
 \qquad\text{for every }m,
 \qquad
 \bs(F_m,0)\longrightarrow\infty.
\]
\end{corollary}

\begin{proof}
Fix $d=9$, choose an integer $M\geq9^{10}$, and choose any admissible
labelling. Let $f$ be the function constructed in the proof of
Theorem~\ref{thm:main} for these fixed parameters. Put
\[
 A=22M^{10},\qquad B=M^2k^9.
\]
Here $L=11$, so~\eqref{eq:final-sensitivity} and
\eqref{eq:final-block-sensitivity} give
$s(f)\leq A$ and $\bs(f,0)\geq B$, while $f(0)=0$.
Moreover,
\[
 \frac{B}{A^2}
   =\frac{(k/M^2)^9}{484}
   \geq\frac{2^9}{484}
   =\frac{128}{121}>1.
\]
In particular, $A>1$, $B>A^2$, and the fixed number
\[
 \alpha=\frac{\log B}{\log A}
\]
is greater than two.

Set $F_1=f$ and, for $m\geq1$, set $F_{m+1}=f\circ F_m$ using disjoint
copies of $F_m$. Inductively, $F_m$ is a total function on $n^m$ bits,
where $n$ is the fixed seed dimension in~\eqref{eq:final-dimension},
and $F_m(0)=0$.
Lemma~\ref{lem:composition} gives
\[
 s(F_m)\leq A^m,
 \qquad
 \bs(F_m,0)\geq B^m.
\]
The latter bound proves nonconstancy and tends to infinity. Since
$A^\alpha=B$, we also obtain
\[
 s(F_m)^\alpha\leq A^{m\alpha}=B^m\leq\bs(F_m,0).
\]
The same estimates give the explicit ratio bound
\[
 \frac{\bs(F_m)}{s(F_m)^2}
 \geq\left(\frac{B}{A^2}\right)^m
 \geq\left(\frac{128}{121}\right)^m.\qedhere
\]
\end{proof}
